% ============================================================================
\section{Universal Foundations: Euclidean Ward Identities, Graded BV Algebra, and the Heat-Kernel Resistor Form \texorpdfstring{$\Omega_\Gamma$}{Omega\_Gamma}}
\label{sec:ward-bv-primer}
% ============================================================================

Following the proximity principle, this chapter builds strictly the \textbf{universal foundation} required across every calculation in the monograph---Euclidean Ward--Takahashi identities, graded Batalin--Vilkovisky (BV) abstract indices, Costello's codifferential heat-kernel parametrix $P(\eps,L)$, and the four-stage Gaussian vertex derivation of the master $(H,T)$ electrical-network differential form $\Omega_\Gamma$---so that we enter multi-loop Feynman-diagram integration immediately in Section~\ref{sec:schwinger-periods-k5}. Sector-specific background (off-shell four-dimensional supersymmetric twists, Chern--Simons transgression and exact Wess--Zumino winding $N_W[g]\in\ZZ$, Whitehead cohomology, Vogel's universal plane, Elliott--Safronov--Williams pure-spinor taxonomy, eleven-dimensional supergravity, Calabi--Yau BCOV polyvectors, and $\fpsl(3|3)$) is placed as dedicated in-situ \textbf{Background Cards} directly alongside the calculations in Sections~\ref{sec:forests-quartic-jets}--\ref{sec:sugra-bcov-holo} where each structure first enters.

\subsection{Euclidean Ward Identities, Graded BV Abstract Indices, and the Four-Stage Heat-Kernel Derivation of \texorpdfstring{$\Omega_\Gamma$}{Omega\_Gamma}}
\label{subsec:bv-graded-module}

In Euclidean quantum field theory with classical action $S_{E,\mathrm{cl}}[\Phi]$ (path-integral Boltzmann weight $\ee^{-S_{E,\mathrm{cl}}/\hbar}$, related to Lorentzian $S_L$ by Wick continuation $S_E = -\ii S_L$), the physical content of a continuous symmetry $\delta\Phi^\Lambda = R^\Lambda[\Phi]$ is its family of \textbf{Ward--Takahashi} (or \textbf{Slavnov--Taylor}) \textbf{identities} generated by functional Schwinger--Dyson integration by parts $\int\mathcal{D}\Phi\,\frac{\delta_l}{\delta\Phi^\Lambda}\bigl(R^\Lambda\mathcal{O}_{\mathrm{tot}}\,\ee^{-S_{E,\mathrm{cl}}/\hbar}\bigr) = 0$. Introducing an ultraviolet cutoff $\eps > 0$ and renormalizing local counterterms deforms both the effective action and the symmetry variation of composite operators into formal loop series $Q_\hbar = Q_0 + \sum_{\substack{n \ge 1\\\mathrm{loop\ orders}}} \hbar^n Q_n$. Introducing \emph{antifields} $\Phi^\ddagger_\Lambda(x)$---auxiliary sources of opposite statistics $|\Phi^\ddagger_\Lambda| \equiv |\Phi^\Lambda| + 1 \pmod 2$ and ghost number $\mathrm{gh}(\Phi^\ddagger_\Lambda) = -1 - \mathrm{gh}(\Phi^\Lambda)$ placed \emph{at the back (on the right)} of symmetry variations inside $S_E[\Phi,\Phi^\ddagger] = S_{E,\mathrm{cl}}[\Phi] + \sum_{\substack{\Lambda\\\mathrm{fields}}}\int\dd^dx\,(s\Phi^\Lambda)\Phi^\ddagger_\Lambda + \mathcal{O}((\Phi^\ddagger)^2)$---and restricting the path integral onto a gauge-fixing \emph{Lagrangian submanifold} $\mathcal{L}_\Psi := \{\Phi^\ddagger_\Lambda = \delta_l\Psi/\delta\Phi^\Lambda\}$ (a maximal isotropic slice inside the field--antifield phase space) packages every Ward identity and anomaly obstruction into a single master differential operator (Table~\ref{tab:correlator-vs-bv}).

\begin{table}[H]
\centering
\scriptsize
\setlength{\tabcolsep}{5pt}
\renewcommand{\arraystretch}{1.20}
\begin{tabularx}{\textwidth}{@{}l>{\raggedright\arraybackslash}X>{\raggedright\arraybackslash}X@{}}
\toprule
\textbf{Physical Concept} & \textbf{Standard Correlator Language} & \textbf{Euclidean BV Master-Action Language} \\
\midrule
Fundamental identity & $\int\mathcal{D}\Phi\,\frac{\delta_l}{\delta\Phi^\Lambda}\bigl(R^\Lambda\mathcal{O}_{\mathrm{tot}}\,\ee^{-S_{E,\mathrm{cl}}/\hbar}\bigr) = 0$ & $-\hbar\int_{\mathcal{L}_\Psi}\mathcal{D}\Phi\;\Delta_{\mathrm{BV}}\bigl(\mathcal{O}_{\mathrm{tot}}\,\ee^{-S_E/\hbar}\bigr) = 0$ \\
Expanded Ward / SD eq. & $\langle\delta\mathcal{O}_{\mathrm{tot}}\rangle - \frac{1}{\hbar}\int\dd^dx\,\epsilon\langle(\partial_\mu j^\mu_E)\mathcal{O}_{\mathrm{tot}}\rangle = 0$ & $\langle(s-\hbar\Delta_{\mathrm{BV}})\mathcal{O}_{\mathrm{tot}}\rangle = \frac{(-1)^{|\mathcal{O}_{\mathrm{tot}}|}}{\hbar}\langle\mathcal{O}_{\mathrm{tot}}(\frac{1}{2}\{S_E,S_E\}-\hbar\Delta_{\mathrm{BV}}S_E)\rangle$ \\
Operator variation & Infinitesimal symmetry variation $\delta\mathcal{O}_r$ & Canonical antibracket $s\mathcal{O}_r := \{S_E, \mathcal{O}_r\}$ \\
Symmetry / anomaly defect & Divergence $+\text{ measure Jacobian }\partial_\mu j^\mu_E$ & Infrared-cutoff QME obstruction $O_\hbar = \frac{1}{2}\{S_E,S_E\}_L - \hbar\Delta_L S_E$ \\
\bottomrule
\end{tabularx}
\caption{Dictionary between Euclidean path-integral Schwinger--Dyson / Ward identities and the BV master identity.}
\label{tab:correlator-vs-bv}
\end{table}

\begin{zkmodule}{Graded BV Abstract Indices, Euclidean Master Equation, and $D^*$ Heat-Kernel Parametrix}{bv-linfty-physics-grading}
Let $\mathcal{E}_{\mathrm{BV}} = \bigoplus_{\substack{g\in\ZZ\\\mathrm{gh}}}\mathcal{E}_{\mathrm{BV}}^g$ be a $\ZZ$-graded vector space (equipped in half-integer spin theories with an independent $\ZZ/2$ parity $|\cdot| \in \ZZ/2$). Using Gothic abstract slots $\mathfrak{A},\mathfrak{B}$ on $\mathcal{E}_{\mathrm{BV}}$ and $\mathcal{E}_{\mathrm{BV}}^*$, pick homogeneous basis vectors $(e_r)^{\mathfrak{A}} \in \mathcal{E}_{\mathrm{BV}}$ of degree $|r|$ and dual coordinate covectors $\Phi^r := (e^r)_{\mathfrak{A}} \in \mathcal{E}_{\mathrm{BV}}^*$ ($\deg\Phi^r = -|r|$): a field configuration $\boldsymbol{\varphi}^{\mathfrak{A}} = \Phi^r(\boldsymbol{\varphi})(e_r)^{\mathfrak{A}}$ pairs vectors $(e_r)^{\mathfrak{A}}$ with coordinate functions $\Phi^r$ on $\mathcal{E}_{\mathrm{BV}}$. Grouping $\Phi^r$ into Darboux pairs $(\phi^k, \phi^\ddagger_k)$ ($|\phi^\ddagger_k| \equiv |\phi^k| + 1 \pmod 2$, $\mathrm{gh}(\phi^\ddagger_k) = -1 - \mathrm{gh}(\phi^k)$), the degree-$(-1)$ symplectic form $\omega_{\mathrm{BV},\mathfrak{A}\mathfrak{B}}$ and its inverse $(\omega_{\mathrm{BV}}^{-1})^{\mathfrak{A}\mathfrak{B}}$ generate the canonical \emph{antibracket} and odd BV Laplacian:
\begin{equation}
\label{eq:bv-locked-conventions}
\{\phi^\ddagger_r, \phi^k\} = \delta_r{}^k,\qquad
\Delta_{\mathrm{BV}} := \sum_{\substack{k\\\mathrm{fields}}}(-1)^{|\phi^k|+1}\frac{\partial_l}{\partial\phi^k}\frac{\partial_l}{\partial\phi^\ddagger_k},\qquad
\Delta_{\mathrm{BV}}^2 = 0.
\end{equation}
With Euclidean path-integral weight $\ee^{-S_E/\hbar}$ ($s\mathcal{O} := \{S_E,\mathcal{O}\}$), the Euclidean Quantum Master Equation reads
\begin{equation}
\label{eq:qme-locked-euclidean}
\tfrac{1}{2}\{S_E, S_E\} - \hbar\,\Delta_{\mathrm{BV}}S_E = 0
\quad\Longleftrightarrow\quad
\bigl(Q_{\mathrm{free}} + \{S_{E,\mathrm{int}}, -\} - \hbar\,\Delta_{\mathrm{BV}}\bigr)\mathbf{1} = 0\quad \bigl(\text{when }\Delta_{\mathrm{BV}}S_{E,\mathrm{free}} = 0\bigr).
\end{equation}

\paragraph{Resolution of the $L_\infty$ target puzzle $\ell_k : (\mathfrak{L}^1)^{\otimes k} \to \mathfrak{L}^2$.}
In homological field-value grading where ghosts sit in $\mathcal{E}_{\mathrm{hom}}^{-1}$, physical fields sit in $\mathcal{E}_{\mathrm{hom}}^0$, and antifields sit in $\mathcal{E}_{\mathrm{hom}}^1$, passing to the \emph{physics-graded $L_\infty$ algebra} $\mathfrak{L}^\bullet := \mathcal{E}_{\mathrm{hom}}[1]$ shifts degrees by $\mathfrak{L}^p := \mathcal{E}_{\mathrm{hom}}^{p-1}$: gauge ghosts sit in $\mathfrak{L}^0$, physical fields sit in $\mathfrak{L}^1$, and antifield/equation-of-motion fibres sit in $\mathfrak{L}^2$. Expanding $S_E(\boldsymbol{\varphi}) = \sum_{\substack{k\ge 1\\\mathrm{arity}}}\frac{1}{(k+1)!}\omega_{\mathrm{BV},\mathfrak{A}\mathfrak{B}}\boldsymbol{\varphi}^{\mathfrak{A}}\ell_k(\boldsymbol{\varphi},\dots,\boldsymbol{\varphi})^{\mathfrak{B}}$ determines the degree-$1$ multilinear brackets $\ell_k : (\mathfrak{L}^1)^{\otimes k} \to \mathfrak{L}^2$ via the symplectic inverse $(\omega_{\mathrm{BV}}^{-1})^{\mathfrak{A}\mathfrak{B}}$, mapping a functional derivative $\delta S_E/\delta\phi^m \in (\mathfrak{L}^1)^*$ into an \emph{antifield-sector vector} $(e^\ddagger)^m \in \mathfrak{L}^2$.

\paragraph{The BV gauge-fixing codifferential $D^*$ and Costello's heat-kernel parametrix.}
On $X = \RR^T_\tau \times \C^H_z$, the free kinetic operator of a holomorphic--topological superfield theory is the mixed de~Rham--Dolbeault differential $D_{\mathrm{free}} := \dd_\tau + \dbar = \dd\tau^\rho\partial_{\tau^\rho} + \dd\bar{z}^{\bar{\jmath}}\dbar_{\bar{\jmath}}$ ($D_{\mathrm{free}}^2 = 0$). Just as Maxwell's kinetic operator $\partial^2 g^{\mu\nu} - \partial^\mu\partial^\nu$ annihilates longitudinal pure-gauge modes until Lorenz gauge $\partial^\mu V_\mu = 0$ is imposed, $D_{\mathrm{free}}$ annihilates every $D_{\mathrm{free}}$-exact shift and lacks an inverse on $\Omega^\bullet(\RR^T)\widehat{\otimes}\Omega^{0,\bullet}(\C^H)$ prior to gauge fixing. For each coordinate vector field $\partial_{\tau^\rho}$ or $\partial_{\bar{z}^{\bar{\jmath}}}$, let $\iota_{\partial_{\tau^\rho}}$ and $\iota_{\partial_{\bar{z}^{\bar{\jmath}}}}$ denote \emph{interior contraction} on differential forms---the odd degree-$(-1)$ derivations defined on $1$-forms by $\iota_{\partial_{\tau^\rho}}(\dd\tau^\sigma) = \delta_\rho{}^\sigma$ and $\iota_{\partial_{\bar{z}^{\bar{\jmath}}}}(\dd\bar{z}^{\bar{\imath}}) = \delta_{\bar{\jmath}}{}^{\bar{\imath}}$ (vanishing on complementary $1$-forms). Using the flat Euclidean metric on $\RR^T\times\C^H$, Costello's \emph{BV gauge-fixing operator} $\QGF = D^*$ is the formal $L^2$-adjoint (Hodge codifferential) of $D_{\mathrm{free}}$, lowering form/ghost degree by $1$:
\begin{equation}
\label{eq:dstar-definition}
D^* := -\delta^{\rho\sigma}\iota_{\partial_{\tau^\rho}}\partial_{\tau^\sigma} - 4\delta^{\bar{\jmath}i}\iota_{\partial_{\bar{z}^{\bar{\jmath}}}}\partial_i,\qquad
(D^*)^2 = 0,\qquad
[D_{\mathrm{free}}, D^*]_+ := D_{\mathrm{free}}D^* + D^* D_{\mathrm{free}} = -\Delta_X,
\end{equation}
where $-\Delta_X := -\delta^{\rho\sigma}\partial_{\tau^\rho}\partial_{\tau^\sigma} - 4\delta^{\bar{\jmath}i}\partial_i\dbar_{\bar{\jmath}}$ is the positive-definite scalar Laplacian (via Cartan's anticommutators $[\dd\tau^\rho,\iota_{\partial_{\tau^\sigma}}]_+ = \delta^\rho{}_\sigma$ and $[\dd\bar{z}^{\bar{\imath}},\iota_{\partial_{\bar{z}^{\bar{\jmath}}}}]_+ = \delta^{\bar{\imath}}{}_{\bar{\jmath}}$). On a $1$-form potential $\mathcal{A}^{(1)} = A_\rho\,\dd\tau^\rho + A_{\bar{\jmath}}\,\dd\bar{z}^{\bar{\jmath}}$, imposing $D^*\mathcal{A}^{(1)} = -\delta^{\rho\sigma}\partial_{\tau^\rho}A_\sigma - 4\delta^{\bar{\jmath}i}\partial_i A_{\bar{\jmath}} = 0$ is literally the differential-form \emph{Lorenz / $\dbar^*$-harmonic gauge condition}, cutting out the gauge-fixed Lagrangian subspace $\ker D^*$. Representing $(-\Delta_X)^{-1} = \int_0^\infty\dd t\,\ee^{t\Delta_X}$ via the heat kernel $K_t = \ee^{t\Delta_X}$ ($\partial_t K_t = \Delta_X K_t = -[D_{\mathrm{free}}, D^*]_+ K_t$) and applying $D^*$ constructs Costello's \emph{parametrix}---a smooth-cutoff propagator $P(\eps,L) := \int_\eps^L\dd t\,D^* K_t$ interpolating between ultraviolet cutoff $\eps > 0$ and infrared scale $L > 0$~\cite{Costello}, obeying $D^* P(\eps,L) = 0$ and $[D_{\mathrm{free}}, P(\eps,L)]_+ = K_\eps - K_L$---with explicit $(T+H)$-form kernel (fixing the factor-$4$ normalization against~\cite{WW24}, Appendix~\ref{app:periods-errata}):
\begin{equation}
\label{eq:ht-heat-kernel-explicit}
[K_t(\tau,z)]_\Omega = \frac{\ee^{-\frac{|\tau|^2}{4t} - \frac{|z|^2}{t}}}{(4\pi t)^{\frac{1}{2}T}(\pi t)^H}\,\dd\tau^1\wedge\cdots\wedge\dd\tau^T\wedge\frac{\ii^H}{2^H}\,\dd\bar{z}^{\bar{1}}\wedge\cdots\wedge\dd\bar{z}^{\bar{H}}.
\end{equation}
\end{zkmodule}

\begin{example}[{Scalar $\phi^4$ vs.\ Holomorphic Chern--Simons in $\mathfrak{L}^\bullet = \mathcal{E}_{\mathrm{hom}}[1]$}]
\label{ex:bv-phi4-hcs}
For Euclidean $\phi^4$ on $\RR^4$, $S_E[\phi,\phi^\ddagger] = \int\dd^4x\,[\frac{1}{2}\phi(-\nabla^2 + m_0^2)\phi + \frac{\lambda_{\mathrm{cpl}}}{4!}\phi^4]$ has $\mathfrak{L}^1 = C^\infty(\RR^4)\cdot e_\phi$ and $\mathfrak{L}^2 = C^\infty(\RR^4)\cdot e_{\phi^\ddagger}$ paired by $\omega_{\mathrm{BV}}(f_1 e_{\phi^\ddagger}, f_2 e_\phi) = \int\dd^4x\,f_1 f_2$. Differentiating $S_E$ with respect to $\phi(x)$ and raising the index by $\omega_{\mathrm{BV}}^{-1}$ lands in $\mathfrak{L}^2$:
\[
\ell_1(\phi\,e_\phi) = \bigl((-\nabla^2 + m_0^2)\phi\bigr)\,e_{\phi^\ddagger} \in \mathfrak{L}^2,\qquad
\ell_3(\phi_1 e_\phi, \phi_2 e_\phi, \phi_3 e_\phi) = \lambda_{\mathrm{cpl}}(\phi_1\phi_2\phi_3)\,e_{\phi^\ddagger} \in \mathfrak{L}^2,
\]
carrying the explicit antifield basis vector $e_{\phi^\ddagger} \in \mathfrak{L}^2$ rather than a bare scalar function! For holomorphic Chern--Simons on $\C^3$ ($\Omega = \dd z^1\wedge\dd z^2\wedge\dd z^3$), $\mathfrak{L}^\bullet = \Omega^{0,\bullet}(\C^3)\otimes\fg$ places ghosts $\ghc^A \in \Omega^{0,0}\otimes\fg = \mathfrak{L}^0$, gauge $(0,1)$-forms $\mathcal{A}^{(0,1)A} \in \mathfrak{L}^1$, antifields $\mathcal{A}^\ddagger \in \Omega^{0,2}\otimes\fg = \mathfrak{L}^2$, and ghost antifields in $\Omega^{0,3}\otimes\fg = \mathfrak{L}^3$, with $\ell_1(\alpha) = \dbar\alpha$ and $\ell_2(\alpha_1,\alpha_2) = [\alpha_1,\alpha_2]$ satisfying graded Jacobi via $\{S_E, S_E\} = 0$.
\end{example}

\begin{lemma}[{Graded Darboux pairing $(\omega_{\mathrm{BV}}^{-1})^{\mathfrak{A}\mathfrak{B}}$, nilpotency $\Delta_{\mathrm{BV}}^2 = 0$, and the Euclidean QME}]
\label{lem:bv-abstract-indices}
\leavevmode
\begin{enumerate}[label=(\arabic*),leftmargin=*]
  \item On a finite-dimensional $\ZZ/2$-graded BV vector space $\mathcal{E}_{\mathrm{BV}}$ with Darboux basis vectors $(e_k)^{\mathfrak{A}}, (e^\ddagger)^{k\mathfrak{A}}$ ($|(e^\ddagger)^k| \equiv |e_k| + 1 \pmod 2$) and coordinate functions $\phi^k := (e^k)_{\mathfrak{A}}, \phi^\ddagger_k := (e^\ddagger_k)_{\mathfrak{A}}$, the odd symplectic inverse $(\omega_{\mathrm{BV}}^{-1})^{(e^\ddagger)^r, e_k} = -\bigl(\omega_{\mathrm{BV}}^{-1}\bigr)^{e_k, (e^\ddagger)^r} = +\delta^r{}_k$ reproduces $\{\phi^\ddagger_r, \phi^k\} = \delta_r{}^k$ and $\{\phi^k, \phi^\ddagger_r\} = -\delta_r{}^k$, and $\Delta_{\mathrm{BV}} = \sum_{\substack{k\\\mathrm{fields}}}(-1)^{|\phi^k|+1}\partial_{\phi^k}\partial_{\phi^\ddagger_k}$ satisfies $\Delta_{\mathrm{BV}}^2 = 0$.

  \item For any Grassmann-even Euclidean master action $S_E$ ($|S_E| = 0$), the graded Leibniz rule $\Delta_{\mathrm{BV}}(F G) = (\Delta_{\mathrm{BV}}F)G + (-1)^{|F|}F(\Delta_{\mathrm{BV}}G) + (-1)^{|F|}\{F, G\}$ implies
  \[
  \Delta_{\mathrm{BV}}\bigl(\ee^{-S_E/\hbar}\bigr) = \frac{1}{\hbar^2}\Bigl(\tfrac{1}{2}\{S_E, S_E\} - \hbar\,\Delta_{\mathrm{BV}}S_E\Bigr)\ee^{-S_E/\hbar}.
  \]
\end{enumerate}
\end{lemma}

\begin{proof}\leavevmode
\begin{proofpart}{Claim 1: Darboux inverse components and $\Delta_{\mathrm{BV}}^2 = 0$}
Setting $\{F, G\} := (F\overleftarrow{\partial}_{\mathfrak{A}})(\omega_{\mathrm{BV}}^{-1})^{\mathfrak{A}\mathfrak{B}}(\overrightarrow{\partial}_{\mathfrak{B}}G)$ on $F = \phi^\ddagger_r$ and $G = \phi^k$ fixes $(\omega_{\mathrm{BV}}^{-1})^{(e^\ddagger)^r, e_k} = +\delta^r{}_k$, while graded odd skew-symmetry $\{G, F\} = -(-1)^{(|F|+1)(|G|+1)}\{F, G\}$ gives $\{\phi^k, \phi^\ddagger_r\} = -(-1)^{(|\phi^k|+1)|\phi^k|}\delta_r{}^k = -\delta_r{}^k$ since $p(p+1)$ is always even. In $\Delta_{\mathrm{BV}}^2 = \sum_{\substack{k,m\\\mathrm{pairs}}}(-1)^{|\phi^k|+|\phi^m|}\partial_{\phi^k}\partial_{\phi^\ddagger_k}\partial_{\phi^m}\partial_{\phi^\ddagger_m}$, each diagonal term $k = m$ squares a Grassmann-odd derivative to zero; for $k \neq m$, commuting the Grassmann-odd operator $\partial_{\phi^m}\partial_{\phi^\ddagger_m}$ (total parity $|\phi^m| + |\phi^m| + 1 \equiv 1 \pmod 2$) past $\partial_{\phi^k}\partial_{\phi^\ddagger_k}$ produces $-1$, cancelling $(k,m)$ against $(m,k)$ pairwise.
\end{proofpart}
\begin{proofpart}{Claim 2: Exponential QME identity}
For even $B := -S_E/\hbar$ ($|B| = 0$), induction on $m \ge 1$ via $\Delta_{\mathrm{BV}}(B^m) = (\Delta_{\mathrm{BV}}B)B^{m-1} + B\Delta_{\mathrm{BV}}(B^{m-1}) + \{B, B^{m-1}\}$ and $\{B, B^{m-1}\} = (m-1)B^{m-2}\{B,B\}$ gives $\Delta_{\mathrm{BV}}(B^m) = m(\Delta_{\mathrm{BV}}B)B^{m-1} + \binom{m}{2}\{B,B\}B^{m-2}$. Dividing by $m!$ and summing over $m \ge 0$ yields $\Delta_{\mathrm{BV}}(\ee^B) = (\Delta_{\mathrm{BV}}B + \frac{1}{2}\{B,B\})\ee^B = (-\hbar^{-1}\Delta_{\mathrm{BV}}S_E + \frac{1}{2}\hbar^{-2}\{S_E,S_E\})\ee^{-S_E/\hbar}$.
\end{proofpart}
\end{proof}

\begin{theorem}[{Four-stage heat-kernel Gaussian derivation of the master $(H,T)$ resistor form $\Omega_\Gamma$~{\cite{Mathews,WW24,Repo9e833c8}}}]
\label{thm:heat-kernel-omega-gamma}
Let $\Gamma = (V,I)$ be a connected oriented $n$-loop multigraph ($n = |I|-|V|+1 \ge 1$) on $X = \RR^T_\tau\times\C^H_z$ ($D := H+T \ge 1$) with grounded root $v_0 \in V$, incidence matrix $B^u{}_e$, cycle matrix $C_r{}^e$, conductance Laplacian $\Lambda^{uv}(t) = B^u{}_e G^{ef}(t)B^v{}_f$, Symanzik cycle-resistance matrix $M_{rs}(t) = C_r{}^e R_{ef}(t)C_s{}^f$ ($U_\Gamma(t) = \det M(t)$), and conserved external holomorphic momenta $\lambda^v{}_{\mathsf{a}}$ ($\sum_{\substack{v\in V\\\mathrm{vertices}}}\lambda^v{}_{\mathsf{a}} = 0$, zero topological external momenta). Integrating the product of extended Schwinger edge forms $\bigwedge_{e\in I}\bigl([K_{t_e}(\tau_e,z_e)]_\Omega - \dd t_e\wedge[D^* K_{t_e}(\tau_e,z_e)]_\Omega\bigr)$ against $\exp\bigl(-\sum_{\substack{u\neq v_0\\\mathrm{ungrounded}}}\lambda^u{}_i z_u{}^i\bigr)$ over the $|V|-1$ ungrounded vertex positions $(\tau_u, z_u, \bar{z}_u) \in X^{|V|-1}$ produces up to canonical orientation sign the projective $(|I|-1)$-form
\begin{equation}
\label{eq:sec3-omega-gamma-derived}
\Omega_\Gamma = U_\Gamma(t)^{-H}\,\rho_n(M)^{\wedge T}\wedge\bigwedge_{i=1}^H\bigwedge_{r=1}^n\psi_{r,i},\qquad
\rho_n(M) := (\det M)^{-1/2}\Pf\!\left(-\tfrac{1}{2}\dd M\,M^{-1}\dd M\right),
\end{equation}
where $\psi_{r,i} = C_r{}^e y^f{}_i(t;\lambda)\delta_{ef}\,\dd t^e$~\eqref{eq:psi-cycle-def} and $\iota_E\Omega_\Gamma = 0$ along $E = \sum_{\substack{e\in I\\\mathrm{edges}}}t^e\partial_{t^e}$.
\end{theorem}

\begin{proof}\leavevmode
\begin{proofpart}{Stage 1: Extended Schwinger edge superform}
Applying $D^*$~\eqref{eq:dstar-definition} to $[K_{t_e}(\tau_e,z_e)]_\Omega$~\eqref{eq:ht-heat-kernel-explicit} uses $-\partial_{\tau^\rho}(\ee^{-|\tau_e|^2/(4t_e)}) = \frac{\tau_{e,\rho}}{2t_e}\ee^{-|\tau_e|^2/(4t_e)}$ and $-4\partial_{z^i}(\ee^{-|z_e|^2/t_e}) = \frac{4\bar{z}_{e,i}}{t_e}\ee^{-|z_e|^2/t_e}$. Combining $[K_{t_e}]_\Omega - \dd t_e\wedge[D^* K_{t_e}]_\Omega$ replaces each coordinate differential in~\eqref{eq:ht-heat-kernel-explicit} by its scale-covariant drift combination $\widehat{\dd\tau}_e^\rho := \dd\tau_e^\rho - \frac{\tau_e^\rho}{2t_e}\dd t_e = t_e^{1/2}\dd(\tau_e^\rho t_e^{-1/2})$ and $\widehat{\dd\bar{z}}_e^{\bar{\imath}} := \dd\bar{z}_e^{\bar{\imath}} - \frac{\bar{z}_e^{\bar{\imath}}}{t_e}\dd t_e = t_e\,\dd(\bar{z}_e^{\bar{\imath}}t_e^{-1})$.
\end{proofpart}
\begin{proofpart}{Stage 2: Holomorphic saddle, Ohm currents, and Kirchhoff determinant $U_\Gamma^{-H}$}
Write edge differences as $z_e{}^i = B^u{}_e z_u{}^i$ and $\tau_e^\rho = B^u{}_e\tau_u^\rho$. In the holomorphic exponent $-\sum_{\substack{e\in I\\\mathrm{edges}}}\frac{|z_e|^2}{t_e} - \sum_{\substack{u\neq v_0\\\mathrm{verts}}}\lambda^u{}_i z_u{}^i = -\bar{z}_u{}^{\bar{\imath}}\delta_{\bar{\imath}i}\Lambda^{uv}z_v{}^i - \lambda^u{}_i z_u{}^i$, integrating $(z_u,\bar{z}_u) \in (\C^H)^{|V|-1}$ fixes the antiholomorphic saddle at $\bar{z}_{u,i}^* = -(\Lambda^{-1})_{uv}\lambda^v{}_i$, so $\frac{\bar{z}_{e,i}^*}{t_e} = -G^{ef}B^u{}_f(\Lambda^{-1})_{uv}\lambda^v{}_i = -y^e{}_i(t;\lambda)$ is minus the \emph{Ohm edge current}! Gaussian integration over the $(|V|-1)H$ complex vertex coordinates cancels $|V|-1$ of the $|I|$ antiholomorphic differentials $\dd\bar{z}_e^{\bar{\imath}}$ and produces $(\det\Lambda)^{-H}\prod_{\substack{e\in I\\\mathrm{edges}}}t_e^{-H} = U_\Gamma(t)^{-H}$ by Kirchhoff's tree--cotree identity $U_\Gamma(t) = (\prod_{\substack{e\in I\\\mathrm{edges}}}t_e)\det\Lambda(t)$ (Proposition~\ref{prop:kirchhoff-duality}(2)), while the $n = |I|-(|V|-1)$ surviving cycle projections of $-\frac{\bar{z}_{e,i}^*}{t_e}\dd t_e = y^e{}_i\,\dd t_e$ assemble into $\bigwedge_{r=1}^n\psi_{r,i}$ for each $i \in \{1,\dots,H\}$.
\end{proofpart}
\begin{proofpart}{Stage 3: Topological Wick Pfaffians $\rho_n(M)^{\wedge T}$ and radial projectivity}
For each topological direction $\rho \in \{1,\dots,T\}$, orthogonal Hodge decomposition $\RR^{|I|} = \im(B^{\mathsf{T}})\oplus_G \im(R^{-1}C^{\mathsf{T}})$ splits $\bigwedge_{e\in I}\widehat{\dd\tau}_e^\rho$ into the $|V|-1$ vertex differentials $\bigwedge_{u\neq v_0}\dd\tau_u^\rho$ wedged with the $n$ cycle $1$-forms $\beta_r^\rho := -\sum_{\substack{e\in I\\\mathrm{edges}}}C_r{}^e\frac{\tau_e^\rho}{2t_e}\dd t_e$ ($\tau_e^\rho = B^u{}_e\tau_u^\rho$), multiplied by the change-of-basis Jacobian $(\det\Lambda)^{1/2}(\det M)^{-1/2}\prod_{\substack{e\in I\\\mathrm{edges}}}t_e^{1/2}$. Integrating $\tau_u^\rho \in \RR^{|V|-1}$ against $\prod_{\substack{e\in I\\\mathrm{edges}}}(4\pi t_e)^{-1/2}\ee^{-\frac{1}{4}\tau_u\Lambda^{uv}\tau_v}$ cancels $(\det\Lambda)^{1/2}\prod_{\substack{e\in I\\\mathrm{edges}}}t_e^{1/2}$, leaves $(\det M)^{-1/2}$, and contracts the zero-mean Gaussian $1$-forms $\beta_r^\rho$ in pairs (so odd $n$ vanishes identically). Using $\langle\tau_u^\rho\tau_v^\rho\rangle = 2(\Lambda^{-1})_{uv}$ and $G - G B^{\mathsf{T}}\Lambda^{-1}B G = C^{\mathsf{T}}M^{-1}C$ on $\dd t_e\wedge\dd t_f$, each Wick pair evaluates to $\langle\beta_r^\rho\wedge\beta_s^\rho\rangle = -\frac{1}{2}\dd M_{rp}(M^{-1})^{pq}\dd M_{qs}$, summing by Wick's Pfaffian theorem to $\Pf(-\frac{1}{2}\dd M\,M^{-1}\dd M)$ and hence $\rho_n(M)$ per coordinate $\rho = 1,\dots,T$. Finally, $\iota_E\psi_{r,i} = C_r{}^e R_{ef}y^f{}_i = 0$ and $\iota_E(\dd M_{rp}(M^{-1})^{pq}\dd M_{qs}) = M_{rp}(M^{-1})^{pq}\dd M_{qs} - \dd M_{rp}(M^{-1})^{pq}M_{qs} = 0$ (since $\iota_E\dd M = M$), proving $\iota_E\rho_n(M) = 0$ and $\iota_E\Omega_\Gamma = 0$.
\end{proofpart}
\end{proof}
